\scanpage{4}
\mspage{- 7 -}
\begin{flushright}
n\textsuperscript{o}\handwritten{354}
\end{flushright}

\uline{est un sous-module gradué} de $M$.

COROLLAIRE 1.- \uline{Si} $N$ \uline{est un sous-module gradué de} $M$ \uline{et si} $(x_{\handwritten{\unsure{\iota|\tau}}})$ \uline{est un système de générateurs de} $N$, \uline{les composantes homogènes des} $x_{\handwritten{\iota}}$ \uline{forment un système de générateurs de} $N$.

COROLLAIRE 2.- \uline{Si} $N$ \uline{est un sous-module gradué de type fini} $M$, $N$ \uline{admet un système fini de générateurs homogènes}.

Il suffit d'appliquer le cor.1 en remarquant qu'un élément de $M$ n'a qu'un nombre fini de composantes homogènes $\ne 0$.

Un sous-module gradué de $A_s$ (resp. $A_d$) est appelé \uline{idéal à gauche} (resp. à \uline{droite}) gradué de l'anneau gradué $A$. Si un sous-anneau $B$ de $A$ est un \uline{sous}-$Z$-module gradué de $A$, on a $(B\cap A_\lambda)(B\cap A_\mu)\subset B\cap A_{\lambda+\mu}$ donc la graduation $(B\cap A_\lambda)$ de $B$ est compatible avec sa structure d'anneau ; on dit alors que $B$ est un \uline{sous-anneau gradué} de $A$.

Il est clair que si $N$ (resp. $B$) est un sous-$A$-module gradué de $M$ (resp. un sous-anneau gradué de $A$), l'injection canonique $N\to M$ (resp. $B\to A$) est un homomorphisme gradué de modules (resp. d'anneaux) de degré 0.

Si $N$ est un sous-module gradué d'un $A$-module gradué $M$, $(M_\lambda)$ la graduation de $M$, les $(M_\lambda+N)/N$, identifiés canoniquement à des sous-modules du module quotient $M/N$, forment une \uline{graduation} compatible avec la structure de ce module pour $\lambda\in\Delta$ ; en effet, ces sous-modules s'identifient aux quotients $M_\lambda/(M_\lambda\cap N)$, et $M/N$ est donc leur somme directe ; en outre, on a $A_\lambda(M_\mu+N)\subset A_\lambda M_\mu+N\subset M_{\lambda+\mu}+N$ donc $A_\lambda((M_\mu+N)/N)\subset(M_{\lambda+\mu}+N)/N$, ce qui établit notre assertion. On dit que la graduation $((M_\lambda+N)/N)$ \uline{est la graduation quotient} de celle de $M$ par $N$, et le module quotient $M/N$, muni de cette graduation, s'appelle le \uline{module quotient gradué} de $M$ par $N$ ; l'homomorphisme canonique $M\to M/N$ est un homomorphisme gradué de degré 0 pour cette graduation.
% archive sheet number: 3

